% Generated from locale/tr/content/proof-theory/normalization/segments.tex; do not hand-edit.


\olfileid[tr]{pt}{nor}{seg}

\olsection{Kesitler ve Kesmeler}

!!^a{introduction} kuralının ardından !!a{elimination} kuralının gelmesi,
bir !!a{proof}{}ın içindeki dolambaçtır; bir !!a{proof}{}ı normalleştirmek de bu tür
dolambaçların tümünü gidermektir. Bununla birlikte \Elim{\lor} ve
\Elim{\lexists} kuralları, ``gidermek istediğimiz dolambaç'' tanımını bile
biraz karmaşıklaştırır. Sorun şudur: Bu kurallar yüzünden bir dolambaçta
!!{elimination} kuralının !!{introduction} kuralını \emph{hemen} izlemesi
garanti değildir. Arada, örneğin bir \Elim{\lor} ya da
\Elim{\lexists} bulunabilir:
\[
\AxiomC{}
\DeduceC{$\lexists[x][!C(x)]$}
\AxiomC{$\Discharge{!A}{x}$}
\DeduceC{$!B$}
\DischargeRule{\Intro{\lif}}{x}
\UnaryInfC{$!A \lif !B$}
\DischargeRule{\Elim{\lexists}}{y}
\BinaryInfC{$!A \lif !B$}
\AxiomC{}
\DeduceC{$!A$}
\RightLabel{\Elim{\lif}}
\BinaryInfC{$!B$}
\DisplayProof
\]
Aslında herhangi sayıda \Elim{\lor} ya da \Elim{\lexists},
\Intro{\lif} ile \Elim{\lif} arasına girebilir; örneğin:
\[
\AxiomC{}
\DeduceC{$\lexists[x][!C(x)]$}
\AxiomC{}
\DeduceC{$!D \lor !E$}
\AxiomC{$\Discharge{!A}{x}$}
\DeduceC{$!B$}
\DischargeRule{\Intro{\lif}}{x}
\UnaryInfC{$!A \lif !B$}
\AxiomC{$\Discharge{!A}{x}$}
\DeduceC{$!B$}
\DischargeRule{\Intro{\lif}}{x}
\UnaryInfC{$!A \lif !B$}
\RightLabel{\Elim{\lor}}
\TrinaryInfC{$!A \lif !B$}
\DischargeRule{\Elim{\lexists}}{y}
\BinaryInfC{$!A \lif !B$}
\AxiomC{}
\DeduceC{$!A$}
\RightLabel{\Elim{\lif}}
\BinaryInfC{$!B$}
\DisplayProof
\]
Bu örnek, birden fazla~\Intro{\lif} sonrasında geldiği için aynı
\Elim{\lif} kuralının birden çok dolambaçta yer alabileceğini de gösterir.
Bütün bu olasılıkları hesaba katmak üzere dolambaçları, bir
\Log{N1i}-ispatındaki belirli türden !!{formula} oluşumu dizilerini
kullanarak tanımlarız; bunlara \emph{kesit} denir. Örneğimizde bu, küçük
bir $\Elim{\lor}$ ya da $\Elim{\lexists}$ öncülünün ardından onun sonucunun
geldiği~$!A \lif !B$ formül oluşumları dizisidir. Örnekte, her biri üç~$!A \lif
!B$ oluşumu içeren böyle iki kesit vardır. Resmî tanım şöyledir.

\begin{defn}
Bir \Log{N1i} !!{proof} örneği $\delta$ içindeki bir \emph{kesit},
$\delta$ içindeki aynı $!A$ !!{formula}{}ünün $!A_1$, \dots, $!A_n$ oluşumlarından
oluşan ve aşağıdaki koşulları sağlayan bir dizidir:
\begin{enumerate}
  \item $!A_1$, \Elim{\lor} ya da \Elim{\lexists} sonucunda elde
  edilmemiştir;
  \item $!A_n$, \Elim{\lor} ya da \Elim{\lexists} kuralının küçük
  öncülü değildir;
  \item $i = 1$, \dots,~$n-1$ için $!A_i$, \Elim{\lor} ya da
  \Elim{\lexists} kuralının küçük öncülü ve $!A_{i+1}$ bu çıkarımın
  sonucudur.
\end{enumerate}
Bu kesitin \emph{uzunluğu}~$n$'dir.
\end{defn}

Her kesit giderme hedefi (yani bir dolambaç) değildir. Hedef olanlara
\emph{kesme} denir.

\begin{defn}
Bir \emph{kesme}, $!A_n$'nin !!a{elimination} çıkarımının büyük öncülü
olduğu ve ayrıca ya $n=1$ olup $!A_1 = !A_n$ bir !!a{introduction}
çıkarımının ya da~\FalseInt'in sonucu olduğu ya da~$n>1$ olduğu bir
kesittir.
\end{defn}

Bir kesme kesitinin !!a{introduction} kuralının sonucuyla başlaması
gerekmediğine dikkat edin. Yalnızca !!a{elimination} kuralının büyük
öncülüyle bitmesi gerekir. Bununla birlikte uzunluğu~$1$ olan bir kesitin,
kesme sayılabilmesi için !!{introduction} kurallarından birinin sonucu
olması gerekir.

\begin{defn}
Bir kesmenin \emph{derecesi}~$\depth{!A}$'dır. !!a{proof}~$\delta$'nın
\emph{kesme derecesi}~$\cutrank{\delta}$,~$\delta$ içindeki kesmelerin en
büyük derecesidir. Bir kesmenin derecesi~$\cutrank{\delta}$ ise, yani
derecesi en büyükse, kesme~$\delta$ içinde \emph{maksimal}dir.
$\delta$'nın \emph{kesme uzunluğu}, bütün maksimal kesmelerinin
uzunlukları toplamıdır.
\end{defn}

